% Slides 1--20
\begin{frame}[plain]
  \begin{tikzpicture}[remember picture,overlay]
    \fill[TgirNavy] (current page.south west) rectangle (current page.north east);
    \fill[TgirTeal] (current page.south west) rectangle ([xshift=10mm]current page.north west);
    \node[anchor=west,text width=.82\paperwidth,align=left] at ([xshift=20mm]current page.west)
    {\usebeamerfont{title}\color{white}\inserttitle\par\vspace{4mm}
     \usebeamerfont{subtitle}\color{TgirMint}\insertsubtitle\par\vspace{7mm}
     \color{white!70}\normalsize\insertauthor\hfill\insertdate};
  \end{tikzpicture}
  \src{Repository history and source tree through local commit f066cce; GitHub repository tglauner/tgir\_quantlib\_tools.}
\end{frame}

\sectionframe{01 / Finance foundations}{A derivative pricer is a disciplined cash-flow machine}{Hull, Options Futures and Other Derivatives; QuantLib documentation.}

\lesson{A price is a conditional present value}
{A derivative is a rule that maps future market states into dated cash flows.}
{Pricing chooses a probability measure and discounts those cash flows with the associated numeraire.}
{Implementation must preserve dates, currencies, signs, conventions, and information available at each decision time.}
{The model is not the product; it is the machinery used to value the product's cash-flow rules.}
{Hull; repository pricing specifications.}

\lesson{Four layers keep pricing logic intelligible}
{Product layer: contractual schedules, coupons, notionals, exercise rights, and settlement rules.}
{Market layer: curves, volatilities, FX spot/forwards, correlations, fixings, and quote conventions.}
{Model and numerical layers: dynamics, calibration, solver, grid, simulation, regression, and tolerances.}
{Most serious pricing defects come from mixing these layers or leaving a convention implicit.}
{ARCHITECTURE.md; docs/latex/documents/callable\_bermudan\_xccy\_swap\_specification\_article.tex.}

\lesson{Discount factors are the common language}
{A discount factor $P(0,T)$ is today's value of one currency unit paid with certainty at $T$.}
{For a deterministic cash flow $C_T$, $PV=C_TP(0,T)$. For stochastic cash flow $X_T$, $PV=\mathbb E^Q[D(0,T)X_T]$.}
{Zero and forward rates are representations derived from the same discount-factor curve.}
{Build and validate discount factors first; every later instrument consumes them.}
{QuantLib term-structure documentation; portfolio.py.}

\lesson{A curve is inferred, not directly observed}
{Markets quote deposits, futures, FRAs, swaps, and OIS rates rather than every maturity's discount factor.}
{Bootstrapping solves instruments sequentially so each quoted instrument reprices on the constructed curve.}
{Interpolation fills unquoted dates; extrapolation policy governs dates beyond the last liquid pillar.}
{Curve construction is already a calibration problem, even before stochastic models appear.}
{portfolio.py build\_sofr\_curve; tests/test\_sofr\_curve.py.}

\lesson{Day counts and calendars move real money}
{Accrual fractions depend on conventions such as Actual/360, Actual/365, or 30/360.}
{Business-day rules move dates around holidays and weekends; spot lags change effective dates.}
{A one-day scheduling mismatch can alter fixings, coupons, exercise eligibility, and discounting.}
{Dates and conventions are model inputs, not cosmetic metadata.}
{QuantLib date/calendar APIs; portfolio.py schedule builders.}

\lesson{An interest-rate swap exchanges coupon systems}
{A payer swap pays fixed and receives floating; a receiver swap does the reverse.}
{Each leg is a signed collection of discounted cash flows with its own schedule and day count.}
{At inception, the par fixed rate makes fixed-leg PV equal floating-leg PV.}
{The swap NPV is linear in cash flows but nonlinear in the curve used to project and discount them.}
{portfolio.py \_make\_ois\_swap; QuantLib swap instruments.}

\begin{frame}{The par swap rate is an annuity-weighted forward rate}
  \[
    K_{\mathrm{par}}=\frac{P(0,T_0)-P(0,T_n)}{\sum_{i=1}^{n}\alpha_iP(0,T_i)},
    \qquad
    A(0)=\sum_{i=1}^{n}\alpha_iP(0,T_i).
  \]
  \begin{itemize}\setlength\itemsep{0.9em}
    \item The numerator represents the floating-leg value in a single-curve idealization.
    \item The denominator $A(0)$ is the fixed-leg annuity, the natural numeraire for swaption formulas.
    \item Multi-curve pricing separates forwarding from discounting but retains the annuity intuition.
  \end{itemize}
  \takeaway{A swap is the bridge from observable term structures to optionality on rates.}
  \src{Standard swap valuation; portfolio.py.}
\end{frame}

\lesson{Risk asks how value changes under controlled perturbations}
{DV01 measures value change for a one-basis-point rate move; bucketed DV01 identifies maturity concentration.}
{Vega measures sensitivity to volatility; delta and gamma describe first- and second-order spot exposure.}
{Bump-and-revalue is simple and general, but every bump must preserve market-data consistency.}
{A reported Greek is incomplete without bump size, units, recalibration rule, and randomness policy.}
{portfolio.py trade\_point\_sensitivities; xccy\_pricer\_diagnostics.py.}

\lesson{Options add nonlinear state dependence}
{A call pays $\max(S_T-K,0)$; a put pays $\max(K-S_T,0)$.}
{Volatility matters because convex payoffs benefit from dispersion even when the forward is unchanged.}
{Exercise style determines when the holder may choose: European once, Bermudan on discrete dates, American continuously.}
{Once choice enters the contract, pricing must represent both uncertainty and optimal decisions.}
{Black and Scholes (1973); Longstaff and Schwartz (2001).}

\lesson{A swaption is an option on a forward swap}
{A payer swaption provides the right to pay fixed; a receiver swaption provides the right to receive fixed.}
{The option's state variable is commonly the forward par swap rate, scaled by the fixed-leg annuity.}
{European swaptions benchmark volatility surfaces and calibrate short-rate models.}
{Bermudan swaptions add an exercise policy across a sequence of remaining swaps.}
{docs/RESEARCH.md; portfolio.py swaption builders.}

\twocolesson{Normal and lognormal volatility answer different questions}
{Bachelier / normal}
{$S_T=S_0+\sigma_N W_T$. Rates may cross zero; volatility is quoted in absolute rate units, often basis points.}
{Black / lognormal}
{$dS_t/S_t=\sigma_LdW_t$. The state stays positive; volatility is quoted as a percentage.}
{The quote convention must match the pricing formula and calibration helper.}
{QuantLib BachelierSwaptionEngine and BlackSwaptionEngine; repository normal-vol matrix.}

\lesson{Calibration translates market prices into model parameters}
{Select liquid instruments whose risks resemble the target product.}
{Minimize a declared error measure: price error, relative price error, or implied-volatility error.}
{Apply bounds, fixed-parameter masks, and convergence criteria to stabilize an underdetermined problem.}
{A calibrated parameter is meaningful only together with the helper set and objective.}
{portfolio.py short-rate calibration; standalone\_xccy\_pricer.py.}

\begin{frame}{RMS error summarizes calibration dispersion}
  \[
  \mathrm{RMS}_{\mathrm{rel}}=
  \sqrt{\frac{1}{N}\sum_{i=1}^{N}
  \left(\frac{V_i^{\mathrm{model}}}{V_i^{\mathrm{market}}}-1\right)^2}.
  \]
  \begin{itemize}\setlength\itemsep{0.8em}
    \item Squaring prevents positive and negative residuals from canceling.
    \item The relative form gives comparable weight across instruments of different price magnitudes.
    \item The repository's XCCY v1 report uses equal helper weights and discloses the exact definition.
  \end{itemize}
  \takeaway{One metric never replaces the instrument-by-instrument residual table.}
  \src{standalone\_xccy\_pricer.py calibrate\_rate\_model; templates/xccy\_callable.html.}
\end{frame}

\lesson{A model is evaluated under a measure}
{The risk-neutral measure is tied to a numeraire such as a domestic money-market account.}
{Changing numeraire changes drifts while leaving correctly priced tradable values invariant.}
{In cross-currency models, the foreign rate drift acquires a quanto correction under the domestic measure.}
{Writing dynamics without the FX quotation and measure is incomplete.}
{docs/latex/documents/callable\_bermudan\_xccy\_quant\_two\_pager\_article.tex.}

\lesson{Trees discretize state; Monte Carlo discretizes paths}
{A tree recombines future states and performs backward induction directly at nodes.}
{Monte Carlo samples full paths and estimates expectations statistically.}
{Trees work especially well for one- or two-factor early-exercise problems; dimensionality grows rapidly.}
{Monte Carlo scales better in factors, but early exercise requires continuation-value estimation.}
{QuantLib TreeSwaptionEngine; Longstaff-Schwartz.}

\lesson{Monte Carlo error declines slowly}
{For $N$ independent paths, the standard error is approximately $s/\sqrt{N}$.}
{Halving error therefore requires roughly four times as many paths.}
{Antithetic variates and common random numbers reduce noise without changing the model.}
{Confidence intervals quantify sampling error, not model risk or exercise-policy bias.}
{standalone\_xccy\_pricer.py \_summary and antithetic normals.}

\lesson{Early exercise is an optimal stopping problem}
{At each exercise date, compare immediate exercise value with conditional continuation value.}
{A tree reads continuation from successor nodes; LSM estimates it by regression on simulated states.}
{The decision must use only information available at that date and must respect event ordering.}
{A policy fitted and valued on the same paths is optimistically biased; split the samples.}
{Longstaff and Schwartz (2001); XCCY pricing specification.}

\lesson{Validation separates implementation from belief}
{Unit tests check deterministic identities and input behavior.}
{Calibration tests confirm quoted instruments reprice within declared tolerances.}
{Martingale, limiting-case, convergence, and independent-benchmark tests interrogate the model.}
{A successful test suite is evidence, not proof that a model is fit for every trading use.}
{tests/; docs/latex/documents/xccy\_callable\_validation\_report\_article.tex.}
